\documentclass[../../main.tex]{subfiles}
\begin{document}

\begin{center}
  {\fontsize{48pt}{58pt}\selectfont \textbf{1993}\normalsize\textbf{年度}}
\end{center}

\begin{table}[htb]
  \centering
  \caption{出題テーマと難易度}
  \begin{tabular}{|c|p{9cm}|c|}
    \hline
    \textbf{問題番号} & \textbf{テーマ} & \textbf{難易度} \\
    \hline
    第1問 & 放物線と直線で囲まれる面積比が一定になる条件 & \\
    \hline
    第2問 & 三角関数の極限と定積分の証明 & \\
    \hline
    第3問 & 4次曲線の接線と交点，線分の長さの最大値 & \\
    \hline
    第4問 & 整数値多項式の証明 & \\
    \hline
    第5問 & 確率（サイコロと内積・鋭角条件） & \\
    \hline
  \end{tabular}
\end{table}

\begin{table}[htb]
  \centering
  \caption{解答}
  \renewcommand{\arraystretch}{1.4}
  \begin{tabular}{|c|c|p{9cm}|}
    \hline
    \textbf{問題番号} & \textbf{小問} & \textbf{答え} \\
    \hline
    第1問 & - & $b=q$ \\
    \hline
    第2問 & (1) & $2n+1$ \\
    \cline{2-3}
    & (2) & 証明問題 \\
    \hline
    第3問 & (1) & $\alpha+\beta=-2t$，$\alpha\beta=3t^2-2a$ \\
    \cline{2-3}
    & (2) & $-\sqrt{\dfrac{a}{3}}<t<\sqrt{\dfrac{a}{3}}$ \\
    \cline{2-3}
    & (3) & $L^2=\{1+16t^2(t^2-a)^2\}(8a-8t^2)$ \\
    \cline{2-3}
    & (4) & 最大値 $\dfrac{4\sqrt{3}}{3}$ \\
    \hline
    第4問 & - & 証明問題 \\
    \hline
    第5問 & - & $\dfrac{605}{1296}$ \\
    \hline
  \end{tabular}
\end{table}

\paragraph{試験形態}
5問．試験時間は要確認．

\end{document}
